\documentclass[12pt]{ctexart}
\usepackage{booktabs,amsmath,graphicx,geometry}
\geometry{a4paper,margin=2.5cm}
\begin{document}
\title{shumo 综合评价报告}
\author{AutoEval-Modeling 自动生成}
\date{2026年05月16日}
\maketitle

\section{模型建立}

本文采用CRITIC法确定指标权重，并运用TOPSIS法进行综合评价。

\subsection{指标权重}
\begin{table}[htbp]\centering
\caption{CRITIC法计算结果}
\begin{tabular}{lcccc}
\toprule
指标 & RD投入_亿元 & 专利授权数_件 & 高新企业数_家 & GDP增长率_百分比 \\
\midrule
权重 & 0.1663 & 0.1642 & 0.1718 & 0.4977 \\
\bottomrule
\end{tabular}\end{table}

\subsection{综合评价结果}
\begin{table}[htbp]\centering
\caption{TOPSIS综合评价结果}
\begin{tabular}{lcc}
\toprule
评价对象 & 综合得分 & 排名 \\
\midrule
成都 & 0.6464 & 1 \\
合肥 & 0.6372 & 2 \\
杭州 & 0.5753 & 3 \\
长沙 & 0.5164 & 4 \\
武汉 & 0.4219 & 5 \\
北京 & 0.3965 & 6 \\
深圳 & 0.3797 & 7 \\
西安 & 0.3329 & 8 \\
上海 & 0.2815 & 9 \\
南京 & 0.1843 & 10 \\
\bottomrule
\end{tabular}\end{table}

\subsection{灵敏度分析}
对各指标权重在 $\pm 20\%$ 范围内进行 OAT 扰动分析，结果表明排名整体部分不稳健。

\end{document}